\documentclass[11pt]{article}
\usepackage{manual_docs_style}
\externaldocument{build/terminology_shared_utilities}[terminology_shared_utilities.pdf]
\externaldocument{build/environment_profiles}[environment_profiles.pdf]
\externaldocument{build/generation_policy_schema}[generation_policy_schema.pdf]

\begin{document}

\docTitle{Run Graph Schema}
\phantomsection\label{docstart:run_graph_schema}

\section*{Overview}

A resolved run graph is the canonical planning artifact for \name{} simulation.
It replaces the old nested \code{model_run_specs.json} topology as the primary source of truth.

The resolved run graph is produced by the candidate-generation stage from:

\begin{DocCode}
models/simulink/<MODEL>/experiment_config.json
models/simulink/<MODEL>/generation_policies/<policy_id>.json
models/simulink/<MODEL>/manual_cases.jsonl          # optional
\end{DocCode}

and written under:

\begin{DocCode}
models/simulink/<MODEL>/plans/<policy_id>/
\end{DocCode}

The core artifacts are:

\begin{DocCode}
run_nodes.jsonl
run_edges.jsonl
generation_report.json
validation_report.json
\end{DocCode}

The graph is called resolved because all random sampling decisions have already been made.
Every run has a stable \code{run_id}, every recipe has a \code{recipe_hash}, and every baseline/intervention/time-zero relationship is explicit.
Stage: Simulate does not sample new runs; it only materializes the exact recipes recorded in \code{run_nodes.jsonl}.

\section*{Artifact Summary}

\begin{center}
\begin{tabular}{p{0.30\linewidth}p{0.60\linewidth}}
\hline
\textbf{Artifact} & \textbf{Purpose} \\
\hline
\code{run_nodes.jsonl} & Flat list of runnable simulation recipes. Each line defines one baseline, intervention, or time-zero/static-change run. \\
\code{run_edges.jsonl} & Graph relationships among runs, such as baseline-to-intervention and intervention-to-time-zero links. \\
\code{generation_report.json} & Generation summary, including skipped parameter-direction requests and coverage statistics. \\
\code{validation_report.json} & Validation summary for schema checks, semantic checks, and policy-level acceptance criteria. \\
\hline
\end{tabular}
\end{center}

\section*{Design Principles}

\begin{itemize}
\item Runs are flat nodes, not nested children.
\item \code{run_id} and \code{recipe_hash} are derived from the canonical recipe.
\item Metadata does not affect run identity.
\item Provenance is recorded as metadata, not as part of the recipe.
\item \code{source} and \code{validation_profile} are separate concepts.
\item The graph contains only runnable runs.
\item Missing or impossible parameter-direction requests are recorded in \code{generation_report.json}, not as null-valued fake runs.
\end{itemize}

\section*{Run Kinds}

The resolved run graph uses three run kinds:

\begin{itemize}
\item \code{baseline}: a no-intervention simulation.
\item \code{intervention}: a simulation where exactly one root-cause parameter changes at an intervention time.
\item \code{time0_baseline}: a static-change probe where the same post-intervention value is active from time zero.
\end{itemize}

The paper describes the third kind as the static-change probe.
The implementation may keep the name \code{time0_baseline} for backward compatibility.

\section*{\texorpdfstring{\titlecode{run_nodes.jsonl}}{run_nodes.jsonl}}

\code{run_nodes.jsonl} is a JSON Lines file.
Each line is one complete runnable run node.

A node has the following top-level fields:

\begin{itemize}
\item \code{run_id}: stable run identifier.
\item \code{kind}: one of \code{baseline}, \code{intervention}, or \code{time0_baseline}.
\item \code{family_id}: identifier for the baseline-centered family this run belongs to.
\item \code{recipe_hash}: hash of the canonical \code{recipe}.
\item \code{recipe}: exact simulator input.
\item \code{metadata}: provenance, policy, generator, and validation-profile information.
\end{itemize}

\subsection*{Minimal Shape}

\begin{DocCode}
{
  "run_id": "run_7fd9daa0055fb3f412ef5e0ab8d4c5ba",
  "kind": "intervention",
  "family_id": "fam_de4d7c0a5335e9005510cc2442d0971e",
  "recipe_hash": "7fd9daa0055fb3f412ef5e0ab8d4c5ba",
  "recipe": {
    "model": "BallDrop",
    "baseline_parameters": {
      "mass": 4.0,
      "drag_coeff": 0.05,
      "restitution": 0.8,
      "gravity": 9.81,
      "initial_height": 3.0,
      "initial_velocity": -2.0
    },
    "intervention": {
      "mode": "at_intervention_time",
      "parameter": "gravity",
      "value": 13.2,
      "time": 4.0
    }
  },
  "metadata": {
    "source": "programmatic",
    "policy_id": "benchmark_v1",
    "validation_profile": "benchmark_candidate",
    "seed": 123,
    "generator": "compile_run_graph.py",
    "generator_version": "1.0.0"
  }
}
\end{DocCode}

\section*{Node Field Definitions}

\subsection*{\code{run_id}}

\code{run_id} is the stable identifier for one runnable recipe.

Recommended form:

\begin{DocCode}
run_<32-hex-recipe-hash>
\end{DocCode}

The \code{run_id} must be unique within a plan.
Preferably, it should be globally stable across plans: if two policies produce the exact same canonical recipe, they should produce the same \code{run_id}.

\subsection*{\code{kind}}

Allowed values are \code{baseline}, \code{intervention}, and \code{time0_baseline}.

\subsection*{\code{family_id}}

\code{family_id} groups one baseline run with its derived intervention and time-zero/static-change runs.
All nodes generated from the same baseline recipe should share the same \code{family_id}.

\subsection*{\code{recipe_hash}}

\code{recipe_hash} is the canonical hash of the \code{recipe} object only.
It must not include \code{metadata}.
Hashing should use normalized JSON values, canonical JSON with sorted keys and compact separators, SHA-256, and the first 32 hexadecimal characters of the digest.

\subsection*{\code{metadata}}

\code{metadata} records provenance and intended use.
It is not part of the recipe hash.
Recommended fields include \code{source}, \code{policy_id}, \code{validation_profile}, \code{seed}, \code{generator}, \code{generator_version}, and optional \code{notes}.

\section*{\texorpdfstring{\titlecode{recipe}}{recipe} Schema}

\subsection*{Baseline Recipe}

A baseline run has \code{intervention.mode = "none"}.

\begin{DocCode}
{
  "model": "BallDrop",
  "baseline_parameters": {
    "mass": 4.0,
    "drag_coeff": 0.05,
    "restitution": 0.8,
    "gravity": 9.81,
    "initial_height": 3.0,
    "initial_velocity": -2.0
  },
  "intervention": {
    "mode": "none",
    "parameter": null,
    "value": null,
    "time": null
  }
}
\end{DocCode}

\subsection*{Intervention Recipe}

An intervention run has \code{intervention.mode = "at_intervention_time"}.

\begin{DocCode}
{
  "model": "BallDrop",
  "baseline_parameters": {
    "mass": 4.0,
    "drag_coeff": 0.05,
    "restitution": 0.8,
    "gravity": 9.81,
    "initial_height": 3.0,
    "initial_velocity": -2.0
  },
  "intervention": {
    "mode": "at_intervention_time",
    "parameter": "gravity",
    "value": 13.2,
    "time": 4.0
  }
}
\end{DocCode}

\subsection*{Time-Zero / Static-Change Recipe}

A time-zero/static-change run has \code{intervention.mode = "from_time_zero"}.
It uses the same changed parameter and post-intervention value as the matched intervention run, but the value is active from time zero.

\begin{DocCode}
{
  "model": "BallDrop",
  "baseline_parameters": {
    "mass": 4.0,
    "drag_coeff": 0.05,
    "restitution": 0.8,
    "gravity": 9.81,
    "initial_height": 3.0,
    "initial_velocity": -2.0
  },
  "intervention": {
    "mode": "from_time_zero",
    "parameter": "gravity",
    "value": 13.2,
    "time": 0.0
  }
}
\end{DocCode}

\section*{\texorpdfstring{\titlecode{baseline_parameters}}{baseline\_parameters}}

\code{baseline_parameters} contains all simulator variables required to initialize the run.
It must include every variable listed in:

\begin{DocCode}
experiment_config.json::exposed_variables.parameters
experiment_config.json::exposed_variables.initial_state
\end{DocCode}

All values must be finite JSON numbers.
Each value must lie inside the corresponding \code{allowed_intervals} from \code{experiment_config.json}.
Variables under \code{parameters} may be intervened on.
Variables under \code{initial_state} are baseline-only variables and must not change after simulation start.

\section*{\texorpdfstring{\titlecode{intervention}}{intervention}}

\code{intervention} defines whether and how the baseline recipe is changed.

Fields:

\begin{itemize}
\item \code{mode}: one of \code{none}, \code{at_intervention_time}, or \code{from_time_zero}.
\item \code{parameter}: changed root-cause parameter, or \code{null} for baseline runs.
\item \code{value}: post-intervention value, or \code{null} for baseline runs.
\item \code{time}: intervention time, \code{0.0} for time-zero/static-change runs, or \code{null} for baseline runs.
\end{itemize}

For \code{at_intervention_time}, \code{time} must satisfy:

\begin{DocCode}
0 < time < experiment_config.json::end_time_input_s
\end{DocCode}

For \code{from_time_zero}, \code{time} should be \code{0.0}.
For \code{none}, \code{parameter}, \code{value}, and \code{time} must be \code{null}.

\section*{\texorpdfstring{\titlecode{run_edges.jsonl}}{run\_edges.jsonl}}

\code{run_edges.jsonl} is a JSON Lines file.
Each line is one relationship between run nodes.
Edges are used to recover the baseline/intervention/time-zero structure without nesting nodes.

\subsection*{Edge Shape}

\begin{DocCode}
{
  "edge_type": "baseline_to_intervention",
  "family_id": "fam_de4d7c0a5335e9005510cc2442d0971e",
  "source_run_id": "run_baseline_...",
  "target_run_id": "run_intervention_...",
  "metadata": {
    "parameter": "gravity",
    "direction": "increase",
    "intervention_time": 4.0
  }
}
\end{DocCode}

Allowed \code{edge_type} values are \code{baseline_to_intervention} and \code{intervention_to_time0_baseline}.
Family membership is primarily represented by \code{family_id}; a separate family-membership edge is not required.

\section*{Provenance and Validation Profiles}

The run graph separates provenance from validation.
\code{source} records who or what created the node, while \code{validation_profile} records what the node is allowed to be used for.
These must not be treated as the same concept.

Allowed \code{source} values are \code{programmatic} and \code{manual}.
Recommended \code{validation_profile} values are \code{diagnostic}, \code{exploratory}, \code{benchmark_candidate}, and \code{paper_selected}.

\section*{Skipped or Impossible Directions}

The run graph contains only runnable run nodes.
It must not contain null-valued fake intervention children.
If a requested parameter-direction pair cannot be sampled, the generator records it in \code{plans/<policy_id>/generation_report.json}.

\begin{DocCode}
{
  "family_id": "fam_de4d7c0a5335e9005510cc2442d0971e",
  "parameter": "restitution",
  "direction": "increase",
  "reason": "baseline value too close to upper allowed bound"
}
\end{DocCode}

Stage: Simulate must ignore skipped directions because they are not run nodes.

\section*{\texorpdfstring{\titlecode{generation_report.json}}{generation\_report.json}}

\code{generation_report.json} records what the candidate-generation stage attempted and produced.
Recommended fields include \code{timestamp}, \code{model}, \code{policy_id}, \code{generator}, \code{generator_version}, \code{seed}, requested and generated family counts, node counts, skipped parameter-directions, and coverage summaries.

\section*{\texorpdfstring{\titlecode{validation_report.json}}{validation\_report.json}}

\code{validation_report.json} records validation results for the resolved run graph.

\begin{DocCode}
{
  "timestamp": "2026-05-21T12:00:00Z",
  "model": "BallDrop",
  "policy_id": "benchmark_v1",
  "status": "pass",
  "checks": {
    "unique_run_ids": "pass",
    "unique_recipe_hashes": "pass",
    "valid_edges": "pass",
    "values_in_allowed_intervals": "pass",
    "interventions_change_one_parameter": "pass",
    "matched_time0_runs": "pass"
  },
  "errors": [],
  "warnings": []
}
\end{DocCode}

Allowed \code{status} values are \code{pass}, \code{pass_with_warnings}, and \code{fail}.
Stage: Simulate may refuse to materialize a plan whose validation report has \code{status = "fail"}.

\section*{Core Validation Rules}

The following rules apply to every resolved run graph:

\begin{itemize}
\item \code{run_nodes.jsonl} must be non-empty.
\item Every line in \code{run_nodes.jsonl} and \code{run_edges.jsonl} must be valid JSON.
\item Every \code{run_id} must be unique.
\item Every node must have \code{run_id}, \code{kind}, \code{family_id}, \code{recipe_hash}, \code{recipe}, and \code{metadata}.
\item Every \code{kind} must be one of \code{baseline}, \code{intervention}, or \code{time0_baseline}.
\item Every recipe must include \code{model}, \code{baseline_parameters}, and \code{intervention}.
\item Every baseline parameter value must be finite.
\item Every baseline recipe must include all exposed parameters and initial-state variables.
\item Every baseline parameter value must respect \code{experiment_config.json::exposed_variables::*::allowed_intervals}.
\item Every intervention run must change exactly one variable listed under \code{experiment_config.json::exposed_variables.parameters}.
\item No intervention may change a variable listed only under \code{experiment_config.json::exposed_variables.initial_state}.
\item Every intervention value must respect the changed parameter's \code{allowed_intervals}.
\item Every \code{at_intervention_time} run must satisfy \code{0 < time < experiment_config.json::end_time_input_s}.
\item Every edge must point to existing run IDs.
\item Every \code{baseline_to_intervention} edge must connect a baseline node to an intervention node.
\item Every \code{intervention_to_time0_baseline} edge must connect an intervention node to a time-zero/static-change node.
\end{itemize}

\section*{Benchmark-Candidate Validation Rules}

The following rules apply to nodes with \code{validation_profile = "benchmark_candidate"} or \code{validation_profile = "paper_selected"}:

\begin{itemize}
\item Every intervention value must differ from the baseline value by at least the configured \termMinAbsDist{}.
\item Every intervention value must differ from the baseline value by at least the configured \termMinSRDDistance{}.
\item Every eligible intervention candidate should have a matched no-intervention baseline run.
\item Every eligible intervention candidate should have a matched time-zero/static-change run.
\item Generated families should provide broad coverage across root-cause parameters.
\item Coverage failures should be recorded in \code{generation_report.json}.
\end{itemize}

Detectability and distinguishability are not guaranteed by the run graph.
They are evaluated after simulation by the metric and oracle stages.

\section*{Runtime Use by Stage: Simulate}

Stage: Simulate reads \code{run_nodes.jsonl} and materializes selected nodes.
For each selected node, it reads the node's \code{recipe}, copies \code{simulink_model_original.mdl} into \code{runs/<run_id>/}, applies \code{recipe.baseline_parameters}, applies \code{recipe.intervention.mode}, serializes the resulting time series to \code{runs/<run_id>/data.parquet}, writes \code{runs/<run_id>/recipe.json}, and updates \code{runs/model_record.json} with \code{run_id}, \code{recipe_hash}, \code{run_type}, and status.

\section*{Relationship to Legacy \texorpdfstring{\titlecode{model_run_specs.json}}{model\_run\_specs.json}}

The old \code{model_run_specs.json} file used a nested structure:

\begin{DocCode}
baseline_uuid -> children -> intervention_uuid -> time0_baseline_uuid
\end{DocCode}

The new run graph replaces this with:

\begin{DocCode}
run_nodes.jsonl  # flat runnable recipes
run_edges.jsonl  # relationships among recipes
\end{DocCode}

\begin{center}
\begin{tabular}{p{0.34\linewidth}p{0.56\linewidth}}
\hline
\textbf{Legacy field} & \textbf{Run-graph replacement} \\
\hline
\code{baseline_uuid} & baseline \code{run_id} \\
\code{intervention_uuid} & intervention \code{run_id} \\
\code{time0_baseline_uuid} & time-zero/static-change \code{run_id} \\
\code{baseline_parameters_hash} & baseline node \code{recipe_hash} \\
\code{parameter_hash} & intervention node \code{recipe_hash} \\
\code{time0_baseline_hash} & time-zero/static-change node \code{recipe_hash} \\
nested \code{children} & \code{run_edges.jsonl} relationships \\
skipped null children & \code{generation_report.json} skipped-direction records \\
\hline
\end{tabular}
\end{center}

\code{model_run_specs.json} may still be generated as a compatibility artifact during migration, but new code should consume the resolved run graph directly.

\section*{Artifact Policy}

\begin{itemize}
\item \code{run_nodes.jsonl} and \code{run_edges.jsonl} are generated artifacts produced by the candidate-generation stage.
\item They should not be manually edited.
\item To change generated candidates, edit the generation policy or manual cases and regenerate the plan.
\item \code{recipe_hash} and \code{run_id} must not depend on metadata.
\item \code{generation_report.json} records what the generator attempted.
\item \code{validation_report.json} records whether the resolved graph is safe to materialize.
\item \code{model_run_specs.json} is legacy compatibility only.
\end{itemize}

\end{document}
